% figures01.tex -- Chapter 1 (INTRODUCTION) figures for Brumfiel, _Geometry_.
% Macro names: \FIG + I (chapter 1) + spelled-out figure number(s), letters only.
% Every constrained point is computed, never eyeballed:
%   bisectors  -> equal-angle construction / (intersection of ...)
%   altitudes  -> ($(A)!(P)!(B)$)
%   midpoints  -> ($(A)!0.5!(B)$)
%   parallels  -> shared direction vector / equal ratios
% Constraint audit: tools/constraints/ch01.py

%--------------------------------------------------------------- Fig 1-1
% Plane figures. Freeform illustrative shapes traced from scan01 p24 (book p.7).
\newcommand{\FIGIONE}{%
\begin{bkfigure}[\linewidth]
\begin{tikzpicture}[scale=0.92, every node/.style={font=\scriptsize}]
  % --- row 1
  \draw[given] (0,0) -- (2.0,0);
  \fill (0,0) circle (1.1pt) (2.0,0) circle (1.1pt);
  \node[vlab, above] at (0,0.02) {$A$};
  \node[vlab, above] at (2.0,0.02) {$B$};
  \node[below=5pt] at (1.0,-0.05) {Line segment};

  \draw[given] (3.3,0) -- (5.3,0);
  \fill (3.3,0) circle (1.1pt);
  \node[vlab, above] at (3.3,0.02) {$O$};
  \node[below=5pt] at (4.3,-0.05) {Ray from $O$};

  \draw[given] (6.6,-0.12) -- (8.7,-0.12);
  \draw[given] (6.6,-0.12) -- (8.25,0.95);
  \node[below=5pt] at (7.6,-0.17) {Angle};

  % --- row 2
  \begin{scope}[yshift=-2.35cm]
    \draw[given] (0.05,0.22) -- (0.95,1.05) -- (1.95,0) -- cycle;
    \node[below=5pt] at (1.0,-0.05) {Triangle};

    \draw[given] (2.85,0) rectangle (4.15,1.3);
    \node[below=5pt] at (3.5,-0.05) {Square};

    \draw[given] (4.95,0) -- (6.75,0) -- (7.35,1.15) -- (5.55,1.15) -- cycle;
    \node[below=5pt] at (6.15,-0.05) {Parallelogram};

    \draw[given] (8.6,0.6) circle (0.62);
    \node[below=5pt] at (8.6,-0.05) {Circle};
  \end{scope}

  % --- row 3
  \begin{scope}[yshift=-4.7cm]
    \draw[given] (0,0) -- (2.0,0) -- (1.65,1.15) -- (0.35,1.15) -- cycle;
    \node[below=5pt] at (1.0,-0.05) {Trapezoid};

    \draw[given] (3.05,0.62) -- (3.85,1.35) -- (5.05,0) -- (3.25,0.28) -- cycle;
    \node[align=center, below=5pt] at (4.0,-0.05)
         {Quadrilateral, or\\4-sided polygon};

    \draw[given] (6.5,0.55) -- (7.05,1.35) -- (8.35,0.62) -- (7.75,0) -- (6.95,0.1) -- cycle;
    \node[align=center, below=5pt] at (7.5,-0.05)
         {Pentagon, or\\5-sided polygon};
  \end{scope}

  % --- row 4
  \begin{scope}[yshift=-7.35cm]
    \draw[given] (3.3,0) rectangle (5.5,1.25);
    \node[below=5pt] at (4.4,-0.05) {Rectangle};
  \end{scope}
\end{tikzpicture}
\figcap{1--1}
\end{bkfigure}}

%--------------------------------------------------------------- Fig 1-2
% Solid figures. Hidden edges dashed, as the book draws them.
\newcommand{\FIGITWO}{%
\begin{bkfigure}[\linewidth]
\begin{tikzpicture}[scale=0.92, every node/.style={font=\scriptsize},
                    hid/.style={fig, dashed, dash pattern=on 2pt off 1.6pt}]
  % ---- cube
  \begin{scope}
    \def\d{0.42}
    \draw[given] (0,0) rectangle (1.25,1.25);
    \draw[given] (\d,\d) -- (\d+1.25,\d) -- (\d+1.25,\d+1.25) -- (\d,\d+1.25);
    \draw[hid]   (\d,\d) -- (\d,\d+1.25);
    \draw[given] (0,1.25) -- (\d,\d+1.25);
    \draw[given] (1.25,1.25) -- (\d+1.25,\d+1.25);
    \draw[given] (1.25,0) -- (\d+1.25,\d);
    \draw[hid]   (0,0) -- (\d,\d);
    \node[below=5pt] at (0.85,-0.05) {Cube};
  \end{scope}
  % ---- box
  \begin{scope}[xshift=3.15cm]
    \def\d{0.42}
    \draw[given] (0,0) rectangle (1.95,1.05);
    \draw[given] (\d,\d) -- (\d+1.95,\d) -- (\d+1.95,\d+1.05) -- (\d,\d+1.05);
    \draw[hid]   (\d,\d) -- (\d,\d+1.05);
    \draw[given] (0,1.05) -- (\d,\d+1.05);
    \draw[given] (1.95,1.05) -- (\d+1.95,\d+1.05);
    \draw[given] (1.95,0) -- (\d+1.95,\d);
    \draw[hid]   (0,0) -- (\d,\d);
    \node[align=center, below=5pt] at (1.2,-0.05)
         {Rectangular\\parallelepiped or box};
  \end{scope}
  % ---- circular cylinder
  \begin{scope}[xshift=6.95cm]
    \draw[given] (0,0) -- (0,1.6);
    \draw[given] (1.3,0) -- (1.3,1.6);
    \draw[given] (0.65,1.6) ellipse (0.65 and 0.24);
    \draw[given] (0.65,0) ++(180:0.65) arc (180:360:0.65 and 0.24);
    \draw[hid]   (0.65,0) ++(180:0.65) arc (180:0:0.65 and 0.24);
    \node[align=center, below=5pt] at (0.65,-0.26) {Circular\\cylinder};
  \end{scope}

  % ---- tetrahedron: as the book draws it, the near vertex is at the front and
  % the single hidden element is the back base edge L--R
  \begin{scope}[yshift=-3.1cm]
    \coordinate (tL) at (0,0.402);
    \coordinate (tR) at (1.517,0.674);
    \coordinate (tF) at (0.817,0);
    \coordinate (tT) at (0.596,1.750);
    \draw[given] (tT) -- (tL) -- (tF) -- (tR) -- cycle;
    \draw[given] (tT) -- (tF);
    \draw[hid]   (tL) -- (tR);
    \node[below=5pt] at (0.76,-0.1) {Tetrahedron};
  \end{scope}
  % ---- pyramid
  \begin{scope}[xshift=3.15cm, yshift=-3.1cm]
    \coordinate (p1) at (0,0.15);
    \coordinate (p2) at (1.45,0);
    \coordinate (p3) at (2.05,0.55);
    \coordinate (p4) at (0.6,0.72);
    \coordinate (pa) at (1.05,1.75);
    \draw[given] (p1) -- (p2) -- (p3);
    \draw[hid]   (p3) -- (p4) -- (p1);
    \draw[given] (p1) -- (pa) -- (p2);
    \draw[given] (p3) -- (pa);
    \draw[hid]   (p4) -- (pa);
    \node[below=5pt] at (1.0,-0.1) {Pyramid};
  \end{scope}
  % ---- prism: upright triangular prism, near vertical edge facing the reader,
  % top face fully visible, only the back bottom edge hidden (as in the book)
  \begin{scope}[xshift=6.95cm, yshift=-3.1cm]
    \coordinate (uL) at (0,1.732);        \coordinate (dL) at (0,0.426);
    \coordinate (uR) at (1.205,1.656);    \coordinate (dR) at (1.205,0.350);
    \coordinate (uF) at (0.266,1.306);    \coordinate (dF) at (0.266,0);
    \draw[given] (uL) -- (uR) -- (uF) -- cycle;   % top face
    \draw[given] (uL) -- (dL);
    \draw[given] (uR) -- (dR);
    \draw[given] (uF) -- (dF);
    \draw[given] (dL) -- (dF) -- (dR);
    \draw[hid]   (dL) -- (dR);
    \node[below=5pt] at (0.60,-0.1) {Prism};
  \end{scope}

  % ---- sphere
  \begin{scope}[yshift=-6.2cm]
    \draw[given] (0.85,0.85) circle (0.78);
    \draw[given] (0.85,0.85) ++(180:0.78) arc (180:360:0.78 and 0.27);
    \draw[hid]   (0.85,0.85) ++(180:0.78) arc (180:0:0.78 and 0.27);
    \node[below=5pt] at (0.85,-0.05) {Sphere};
  \end{scope}
  % ---- torus
  \begin{scope}[xshift=3.15cm, yshift=-6.2cm]
    \draw[given] (1.05,0.85) circle (0.82);
    \draw[given] (1.12,0.92) circle (0.40);
    \draw[given] (0.72,0.92) arc (180:300:0.40 and 0.16);
    \node[align=center, below=5pt] at (1.05,-0.05)
         {Torus, or doughnut,\\or anchor ring};
  \end{scope}
  % ---- cone
  \begin{scope}[xshift=6.95cm, yshift=-6.2cm]
    \draw[given] (0,0.35) -- (0.78,1.85) -- (1.56,0.35);
    \draw[given] (0.78,0.35) ++(180:0.78) arc (180:360:0.78 and 0.27);
    \draw[hid]   (0.78,0.35) ++(180:0.78) arc (180:0:0.78 and 0.27);
    \node[below=5pt] at (0.78,-0.05) {Cone};
  \end{scope}
\end{tikzpicture}
\figcap{1--2}
\end{bkfigure}}

%--------------------------------------------------------------- Fig 1-3
% "Strange" sets of points. ORGANIC -- traced against scan02 p2 (book p.9):
% damped zigzag, chain of shrinking circles, recursively quartered square,
% face built from circles.
\newcommand{\FIGITHREE}{%
\begin{bkfigure}[0.92\linewidth]
\begin{tikzpicture}[scale=1.0]
  % ---- damped zigzag: amplitude falls geometrically left to right
  \begin{scope}[shift={(0,3.05)}]
    % the path must OPEN with a coordinate -- starting it with "--" emits a
    % lineto with no current point and some viewers drop the whole zigzag
    \draw[given] (0,0.42)
      \foreach \i [evaluate={\x=0.30*\i; \s=ifthenelse(mod(\i,2)==0,1,-1);
                            \a=0.42*pow(0.80,\i)} ] in {1,...,15} { -- (\x,{\s*\a}) }
      -- (4.85,0);
  \end{scope}
  % ---- chain of shrinking circles converging to a point
  \begin{scope}[shift={(0.15,1.55)}]
    \foreach \i [evaluate={\r=0.235*pow(0.80,\i); \c=0.44*(1-pow(0.80,\i))/0.20}] in {0,...,9}
      \draw[given] ({\c+\r},0) ellipse ({\r} and {1.28*\r});
    \draw[given] (2.30,0.05) -- (2.62,0) -- (2.30,-0.05);
  \end{scope}
  % ---- square quartered again and again in its upper-right corner
  \begin{scope}[shift={(3.15,0.60)}]
    \def\S{2.0}
    \draw[given] (0,0) rectangle (\S,\S);
    \draw[given] (0,{\S/2}) -- (\S,{\S/2});
    \draw[given] ({\S/2},0) -- ({\S/2},\S);
    \draw[given] ({\S/2},{3*\S/4}) -- (\S,{3*\S/4});
    \draw[given] ({3*\S/4},{\S/2}) -- ({3*\S/4},\S);
    \draw[given] ({3*\S/4},{7*\S/8}) -- (\S,{7*\S/8});
    \draw[given] ({7*\S/8},{3*\S/4}) -- ({7*\S/8},\S);
    \draw[given] ({7*\S/8},{15*\S/16}) -- (\S,{15*\S/16});
    \draw[given] ({15*\S/16},{7*\S/8}) -- ({15*\S/16},\S);
  \end{scope}
  % ---- face made of circles
  \begin{scope}[shift={(1.55,-1.25)}]
    \draw[given] (0,0) circle (0.80 and 0.86);      % head
    \draw[given] (-0.86,0.02) circle (0.10);        % ears
    \draw[given] ( 0.86,0.02) circle (0.10);
    \draw[given] (-0.30,0.26) circle (0.15);        % eyes
    \draw[given] ( 0.30,0.26) circle (0.15);
    \draw[given] (-0.30,0.26) circle (0.06);
    \draw[given] ( 0.30,0.26) circle (0.06);
    \draw[given] (0,-0.06) circle (0.09 and 0.14);  % nose
    \draw[given] (0,-0.48) circle (0.17);           % mouth
  \end{scope}
\end{tikzpicture}
\figcap{1--3}
\end{bkfigure}}

%--------------------------------------------------------------- Fig 1-4
% Angle: vertex, interior, exterior.
\newcommand{\FIGIFOUR}{%
\begin{bkfigure}[0.60\linewidth]
\begin{tikzpicture}[scale=1.0]
  \coordinate (A) at (0,0);
  \draw[given] (A) -- (2.35,1.20);
  \draw[given] (A) -- (2.35,-1.20);
  \node[vlab] at (-0.20,0) {$A$};
  \node[slab] at (-0.42,0.42) {Vertex};
  \node[slab] at (1.02,0) {Interior};
  \node[slab] at (0.62,-0.88) {Exterior};
\end{tikzpicture}
\figcap{1--4}
\end{bkfigure}}

%--------------------------------------------------------------- Figs 1-5, 1-6
% 1-5: ray AC bisects angle BAD -- C placed at exactly half the angle of B.
% 1-6: acute (< 90) and obtuse (> 90) angles.
\newcommand{\FIGIFIVESIX}{%
\begin{bkfigure}[\linewidth]
\begin{minipage}[t]{0.48\linewidth}\centering
\begin{tikzpicture}[scale=1.0]
  \coordinate (A) at (0,0);
  \coordinate (D) at (2.30,0);
  \coordinate (B) at (58:2.30);
  \coordinate (C) at (29:2.15);       % exact bisector of angle BAD
  \draw[given] (A) -- (D);
  \draw[given] (A) -- (B);
  \draw[given] (A) -- (C);
  \node[vlab, left]  at (A) {$A$};
  \node[vlab, right] at (D) {$D$};
  \node[vlab, above] at (B) {$B$};
  \node[vlab, right] at (C) {$C$};
\end{tikzpicture}
\figcap{1--5}
\end{minipage}\hfill
\begin{minipage}[t]{0.48\linewidth}\centering
\begin{tikzpicture}[scale=1.0]
  % acute: vertex bottom-left, 32 degrees
  \draw[given] (0,0) -- (1.55,0);
  \draw[given] (0,0) -- (32:1.55);
  \node[slab, above] at (0.72,0.85) {Acute angle};
  % obtuse: vertex bottom-right; leftward ray at 180 deg, second ray at 47 deg,
  % so the angle at the vertex is 180 - 47 = 133 deg
  \draw[given] (3.35,0) -- (1.95,0);
  \draw[given] (3.35,0) -- ($(3.35,0)+(47:1.55)$);
  \node[slab, above] at (2.85,1.15) {Obtuse angle};
\end{tikzpicture}
\figcap{1--6}
\end{minipage}
\end{bkfigure}}

%--------------------------------------------------------------- Figs 1-7, 1-8
% 1-7: right triangle, right angle at the lower-left vertex.
% 1-8: angle AOB is a right angle; C interior.
\newcommand{\FIGISEVENEIGHT}{%
\begin{bkfigure}[\linewidth]
\begin{minipage}[t]{0.48\linewidth}\centering
\begin{tikzpicture}[scale=1.0]
  \coordinate (R) at (0,0);
  \coordinate (T) at (0,1.75);
  \coordinate (S) at (2.15,0);
  \draw[given] (R) -- (T) -- (S) -- cycle;
  \draw[fig] (0,0.17) -- (0.17,0.17) -- (0.17,0);
  \node[slab] at (0.56,0.30) {90\textdegree};
\end{tikzpicture}
\figcap{1--7}
\end{minipage}\hfill
\begin{minipage}[t]{0.48\linewidth}\centering
\begin{tikzpicture}[scale=1.0]
  \coordinate (O) at (0,0);
  \coordinate (A) at (0,1.75);
  \coordinate (B) at (1.85,0);
  \coordinate (C) at (33:1.62);
  \draw[given] (O) -- (A);
  \draw[given] (O) -- (B);
  \draw[given] (O) -- (C);
  \fill (A) circle (1.2pt) (B) circle (1.2pt) (C) circle (1.2pt);
  \node[vlab, above] at (A) {$A$};
  \node[vlab, right] at (B) {$B$};
  \node[vlab, right] at (C) {$C$};
  \node[vlab, below left, xshift=-1pt, yshift=-0pt] at (O) {$O$};
\end{tikzpicture}
\figcap{1--8}
\end{minipage}
\end{bkfigure}}

%--------------------------------------------------------------- Figs 1-9, 1-10
% 1-9: A, O, C collinear on l (so angle AOB and angle BOC are supplementary).
% 1-10: straight angle A-O-B.
\newcommand{\FIGININETEN}{%
\begin{bkfigure}[\linewidth]
\begin{minipage}[t]{0.48\linewidth}\centering
\begin{tikzpicture}[scale=1.0]
  \coordinate (A) at (-1.55,0);
  \coordinate (O) at (0,0);
  \coordinate (C) at (1.55,0);
  \coordinate (B) at (38:1.62);
  \draw[given] (A) -- (C);
  \draw[given] (O) -- (B);
  \fill (A) circle (1.2pt) (O) circle (1.2pt) (C) circle (1.2pt) (B) circle (1.2pt);
  \node[vlab, left]  at (A) {$A$};
  \node[vlab, right] at (C) {$C$};
  \node[vlab, above right] at (B) {$B$};
  \node[vlab, below] at (O) {$O$};
  \node[vlab, below] at (0.80,-0.02) {$l$};
\end{tikzpicture}
\figcap{1--9}
\end{minipage}\hfill
\begin{minipage}[t]{0.48\linewidth}\centering
\begin{tikzpicture}[scale=1.0]
  \coordinate (A) at (-1.45,0);
  \coordinate (O) at (0,0);
  \coordinate (B) at (1.45,0);
  \draw[given] (A) -- (B);
  \fill (A) circle (1.2pt) (O) circle (1.2pt) (B) circle (1.2pt);
  \node[vlab, above] at (A) {$A$};
  \node[vlab, above] at (B) {$B$};
  \node[vlab, below] at (O) {$O$};
\end{tikzpicture}
\figcap{1--10}
\end{minipage}
\end{bkfigure}}

%--------------------------------------------------------------- Fig 1-11
\newcommand{\FIGIELEVEN}{%
\begin{bkfigure}[0.52\linewidth]
\begin{tikzpicture}[scale=1.0]
  % vertex ratios measured off the book's own quadrilateral
  \draw[given] (0,0.310) -- (0.558,1.565) -- (2.200,0.558) -- (1.270,0) -- cycle;
\end{tikzpicture}
\figcap{1--11}
\end{bkfigure}}

%--------------------------------------------------------------- Figs 1-12, 1-13
% 1-12: equilateral triangle by the two-arc construction (Euclid I.1).
% 1-13: an angle and its copy -- both drawn at exactly 52 degrees.
\newcommand{\FIGITWELVETHIRTEEN}{%
\begin{bkfigure}[\linewidth]
\begin{minipage}[t]{0.46\linewidth}\centering
\begin{tikzpicture}[scale=1.0]
  \def\s{2.05}
  \coordinate (A) at (0,0);
  \coordinate (B) at (\s,0);
  \coordinate (C) at (60:\s);          % exact equilateral apex
  \draw[given] (A) -- (B) -- (C) -- cycle;
  \draw[fig] (A) ++(-18:\s) arc (-18:78:\s);
  \draw[fig] (B) ++(198:\s) arc (198:102:\s);
  \draw[fig] ($(C)+(-0.09,-0.09)$) -- ($(C)+(0.09,0.09)$);
  \draw[fig] ($(C)+(-0.09,0.09)$) -- ($(C)+(0.09,-0.09)$);
  % compass-point marks where each arc passes through the opposite base vertex
  \draw[fig] ($(A)+(135:0.11)$) -- ($(A)+(-45:0.11)$);
  \draw[fig] ($(B)+(45:0.11)$)  -- ($(B)+(225:0.11)$);
\end{tikzpicture}
\figcap{1--12}
\end{minipage}\hfill
\begin{minipage}[t]{0.50\linewidth}\centering
\begin{tikzpicture}[scale=1.0]
  \foreach \x/\lo/\lb/\la in {0/{O}/{B}/{A}, 2.75/{O\mkern2mu'}/{B\mkern2mu'}/{A\mkern2mu'}} {
    \begin{scope}[xshift=\x cm]
      \coordinate (V) at (0,0);
      \draw[given] (V) -- (1.45,0);
      \draw[given] (V) -- (60:1.60);
      \draw[fig] (V) ++(0:0.62) arc (0:60:0.62);
      % the book marks where the arc cuts each ray: a tick on OB, a cross on OA
      \draw[fig] ($(V)+(0.62,-0.10)$) -- ($(V)+(0.62,0.10)$);
      \draw[fig] ($(V)+(60:0.62)+(15:0.11)$)  -- ($(V)+(60:0.62)+(195:0.11)$);
      \draw[fig] ($(V)+(60:0.62)+(105:0.11)$) -- ($(V)+(60:0.62)+(285:0.11)$);
      \node[vlab, below left, xshift=-0pt, yshift=-0pt] at (V) {$\lo$};
      \node[vlab, right] at (1.45,0) {$\lb$};
      \node[vlab, above] at (60:1.60) {$\la$};
    \end{scope}}
\end{tikzpicture}
\figcap{1--13}
\end{minipage}
\end{bkfigure}}

%--------------------------------------------------------------- Fig 1-14
% Compass designs. ORGANIC -- traced against scan02 p9 (book p.14):
% (left) four circles of radius r whose centres lie at distance r from a common
% point; (right) the six-petal rosette inscribed in a circle.
\newcommand{\FIGIFOURTEEN}{%
\begin{bkfigure}[0.86\linewidth]
\begin{tikzpicture}[scale=1.0]
  % four-circle design
  \begin{scope}
    \def\r{0.78}
    \foreach \a in {0,90,180,270} \draw[given] (\a:\r) circle (\r);
  \end{scope}
  % six-petal rosette
  \begin{scope}[xshift=4.15cm]
    \def\R{1.12}
    \begin{scope}
      \clip (0,0) circle (\R);
      \foreach \a in {0,60,...,300} \draw[given] (\a:\R) circle (\R);
    \end{scope}
    \draw[given] (0,0) circle (\R);
  \end{scope}
\end{tikzpicture}
\figcap{1--14}
\end{bkfigure}}

%--------------------------------------------------------------- Figs 1-15, 1-16
% 1-15: angle bisected with compass; Z is the exact arc crossing, so VZ bisects.
% 1-16: adjacent angles AOB and BOC sharing ray OB.
\newcommand{\FIGIFIFTEENSIXTEEN}{%
\begin{bkfigure}[\linewidth]
\begin{minipage}[t]{0.48\linewidth}\centering
\begin{tikzpicture}[scale=1.0]
  \coordinate (V) at (0,0);
  \draw[given] (V) -- (26:2.15);
  \draw[given] (V) -- (-26:2.15);
  \coordinate (X) at (26:0.72);
  \coordinate (Y) at (-26:0.72);
  \draw[fig] (V) ++(34:0.72) arc (34:-34:0.72);
  % arcs of radius 0.716170 about X and Y meet at Z on the bisector; as in the
  % book only their crossing is marked, the two small arcs are not drawn
  \coordinate (Z) at (1.29,0);
  \draw[given] (V) -- (2.15,0);
  \draw[fig] ($(Z)+(-0.09,-0.09)$) -- ($(Z)+(0.09,0.09)$);
  \draw[fig] ($(Z)+(-0.09,0.09)$) -- ($(Z)+(0.09,-0.09)$);
\end{tikzpicture}
\figcap{1--15}
\end{minipage}\hfill
\begin{minipage}[t]{0.48\linewidth}\centering
\begin{tikzpicture}[scale=1.0]
  \coordinate (O) at (0,0);
  \coordinate (A) at (36:2.05);
  \coordinate (B) at (-2:2.05);
  \coordinate (C) at (-24:2.05);
  \draw[given] (O) -- (A);
  \draw[given] (O) -- (B);
  \draw[given] (O) -- (C);
  \node[vlab, left]  at (O) {$O$};
  \node[vlab, above right] at (A) {$A$};
  \node[vlab, right] at (B) {$B$};
  \node[vlab, right] at (C) {$C$};
\end{tikzpicture}
\figcap{1--16}
\end{minipage}
\end{bkfigure}}

%--------------------------------------------------------------- Figs 1-17, 1-18
% 1-17: three pairs of angles, none of them adjacent (traced from scan02 p9):
%   (a) two angles at two DIFFERENT vertices;
%   (b) two angles sharing a vertex and a side but overlapping interiors;
%   (c) two angles sharing only the vertex -- no common side.
% 1-18: copying a segment onto a line at A.
\newcommand{\FIGISEVENTEENEIGHTEEN}{%
\begin{bkfigure}[\linewidth]
\begin{minipage}[t]{0.56\linewidth}\centering
\begin{tikzpicture}[scale=0.95]
  % (a) two nearby but distinct vertices
  \begin{scope}
    \coordinate (V) at (0,0);
    \coordinate (W) at (0.16,-0.14);
    \draw[given] (V) -- (34:1.75);
    \draw[given] (V) -- (-9:1.70);
    \draw[given] (W) -- ($(W)+(-6:1.55)$);
    \draw[given] (W) -- ($(W)+(-44:1.30)$);
    \node[vlab] at (12.5:0.80) {$A$};
    \node[vlab] at ($(W)+(-25:0.98)$) {$B$};
  \end{scope}
  % (b) shared vertex and side, interiors overlap
  \begin{scope}[xshift=3.35cm, yshift=0.05cm]
    \coordinate (V) at (1.30,0);
    \draw[given] (V) -- ($(V)+(140:1.45)$);
    \draw[given] (V) -- ($(V)+(218:1.35)$);
    \draw[given] (V) -- ($(V)+(2:1.15)$);
    \node[vlab] at ($(V)+(35:0.52)$) {$A$};
    \node[vlab] at ($(V)+(-35:0.52)$) {$B$};
  \end{scope}
  % (c) shared vertex only
  \begin{scope}[xshift=1.55cm, yshift=-2.55cm]
    \coordinate (V) at (0,0);
    \draw[given] (V) -- (92:1.55);
    \draw[given] (V) -- (62:1.45);
    \draw[given] (V) -- (-10:1.50);
    \draw[given] (V) -- (-38:1.45);
    \node[vlab] at (77:1.15) {$A$};
    \node[vlab] at (-24:1.25) {$B$};
  \end{scope}
\end{tikzpicture}
\figcap{1--17}
\end{minipage}\hfill
\begin{minipage}[t]{0.40\linewidth}\centering
\begin{tikzpicture}[scale=1.0]
  \draw[given] (0.28,0.52) -- (2.05,0.52);        % the given segment
  \draw[given] (0,0) -- (2.30,0);                  % the line
  \coordinate (A) at (0.28,0);
  \fill (A) circle (1.2pt);
  \draw[fig] (2.05,-0.13) -- (2.05,0.13);          % compass mark at B
  \node[vlab, below left, xshift=2pt, yshift=-0pt] at (A) {$A$};
  \node[vlab, above right, xshift=-3pt, yshift=0pt] at (2.05,0.02) {$B$};
\end{tikzpicture}
\figcap{1--18}
\end{minipage}
\end{bkfigure}}

%--------------------------------------------------------------- Figs 1-19, 1-20
% 1-19: perpendicular bisector of AB; C is the exact midpoint, l is exactly
%       perpendicular (vertical through C).
% 1-20: perpendicular to l at a point O on l.
\newcommand{\FIGININETEENTWENTY}{%
\begin{bkfigure}[\linewidth]
\begin{minipage}[t]{0.48\linewidth}\centering
\begin{tikzpicture}[scale=1.0]
  % drawn 1.4x the first draft's size, to the book's own proportions
  \coordinate (A) at (-1.89,0);
  \coordinate (B) at (1.89,0);
  \coordinate (C) at ($(A)!0.5!(B)$);     % exact midpoint
  \draw[given] (A) -- (B);
  \draw[given] (0,-2.17) -- (0,2.17);
  % compass arcs of radius 2.45 about A and B; the sweep runs past each
  % crossing, as the book draws it. Chained 23-degree pieces keep every
  % Bezier tight against its own arc.
  \draw[aux] (A) ++(-46:2.45) arc (-46:-23:2.45) arc (-23:0:2.45)
                              arc (0:23:2.45)    arc (23:46:2.45);
  \draw[aux] (B) ++(134:2.45) arc (134:157:2.45) arc (157:180:2.45)
                              arc (180:203:2.45) arc (203:226:2.45);
  \foreach \y in {1.558974,-1.558974} {   % sqrt(2.45^2 - 1.89^2)
    \draw[fig] (-0.12,\y-0.12) -- (0.12,\y+0.12);
    \draw[fig] (-0.12,\y+0.12) -- (0.12,\y-0.12);}
  % where the arcs cut AB (x = -1.89+2.45 and 1.89-2.45)
  \foreach \x in {-0.56,0.56} \draw[fig] (\x,-0.15) -- (\x,0.15);
  \fill (A) circle (1.2pt) (B) circle (1.2pt);
  \node[vlab, left]  at (A) {$A$};
  \node[vlab, right] at (B) {$B$};
  \node[vlab] at (0.20,-0.33) {$C$};
  \node[vlab, below right, xshift=0pt, yshift=-1pt] at (0,-2.17) {$l$};
\end{tikzpicture}
\figcap{1--19}
\end{minipage}\hfill
\begin{minipage}[t]{0.48\linewidth}\centering
\begin{tikzpicture}[scale=1.0]
  % same 1.4x enlargement as Fig 1-19, so the pair matches
  \coordinate (O) at (0,0);
  \draw[given] (-2.03,0) -- (2.03,0);
  \draw[given] (0,-2.03) -- (0,2.17);
  \foreach \x in {-1.12,1.12} \draw[fig] (\x,-0.18) -- (\x,0.18);
  \draw[aux] (-1.12,0) ++(-58:1.61)
      arc (-58:-38.667:1.61) arc (-38.667:-19.333:1.61) arc (-19.333:0:1.61)
      arc (0:19.333:1.61)    arc (19.333:38.667:1.61)   arc (38.667:58:1.61);
  \draw[aux] ( 1.12,0) ++(122:1.61)
      arc (122:141.333:1.61)   arc (141.333:160.667:1.61) arc (160.667:180:1.61)
      arc (180:199.333:1.61)   arc (199.333:218.667:1.61) arc (218.667:238:1.61);
  \foreach \y in {1.156590,-1.156590} {   % sqrt(1.61^2 - 1.12^2)
    \draw[fig] (-0.12,\y-0.12) -- (0.12,\y+0.12);
    \draw[fig] (-0.12,\y+0.12) -- (0.12,\y-0.12);}
  \node[vlab, above right, xshift=0pt, yshift=0pt] at (0,2.17) {$m$};
  \node[vlab, left] at (-2.03,0) {$l$};
  \node[vlab] at (0.20,-0.33) {$O$};
\end{tikzpicture}
\figcap{1--20}
\end{minipage}
\end{bkfigure}}

%--------------------------------------------------------------- Figs 1-21, 1-22
% 1-21: perpendicular to l from an external point P (m is exactly vertical).
% 1-22: the three medians; O is the exact centroid.
\newcommand{\FIGITWENTYONETWENTYTWO}{%
\begin{bkfigure}[\linewidth]
\begin{minipage}[t]{0.46\linewidth}\centering
\begin{tikzpicture}[scale=1.0]
  \coordinate (P) at (0,1.30);
  \draw[given] (-1.60,0) -- (1.60,0);
  \draw[given] (0,-1.15) -- (0,1.55);
  % arc of radius 1.75 about P cuts l at x = +- sqrt(1.75^2 - 1.30^2) = 1.171538;
  % the crossings are at 227.95 deg and 312.05 deg, so the sweep runs past them
  \draw[fig] (P) ++(222:1.75) arc (222:318:1.75);
  \draw[fig] (-0.09,1.30-0.09) -- (0.09,1.30+0.09);
  \draw[fig] (-0.09,1.30+0.09) -- (0.09,1.30-0.09);
  % arcs of radius 1.45 about those points meet at y = -sqrt(1.45^2 - 1.171538^2)
  \draw[fig] (-0.09,-0.854400-0.09) -- (0.09,-0.854400+0.09);
  \draw[fig] (-0.09,-0.854400+0.09) -- (0.09,-0.854400-0.09);
  \node[vlab, left] at (P) {$P$};
  \node[vlab, left] at (-1.60,0) {$l$};
  \node[vlab, below] at (0,-1.15) {$m$};
\end{tikzpicture}
\figcap{1--21}
\end{minipage}\hfill
\begin{minipage}[t]{0.50\linewidth}\centering
\begin{tikzpicture}[scale=1.0]
  % proportions from the book: apex at 0.64 of the base, height 0.627 of the base
  \coordinate (A) at (0,0);
  \coordinate (B) at (3.40,0);
  \coordinate (C) at (2.176,2.132);
  \coordinate (Mc) at ($(A)!0.5!(B)$);
  \coordinate (Ma) at ($(B)!0.5!(C)$);
  \coordinate (Mb) at ($(C)!0.5!(A)$);
  \coordinate (O)  at ($(A)!0.6666666!(Ma)$);   % exact centroid
  \draw[given] (A) -- (B) -- (C) -- cycle;
  \draw[aux] (A) -- (Ma);
  \draw[aux] (B) -- (Mb);
  \draw[aux] (C) -- (Mc);
  \draw[fig] ($(O)+(-0.08,-0.08)$) -- ($(O)+(0.08,0.08)$);
  \draw[fig] ($(O)+(-0.08,0.08)$) -- ($(O)+(0.08,-0.08)$);
  % 296.33 deg bisects the widest gap between the six median rays at O
  \node[vlab] at ($(O)+(296.33:0.47)$) {$O$};
\end{tikzpicture}
\figcap{1--22}
\end{minipage}
\end{bkfigure}}

%--------------------------------------------------------------- Fig 1-23
% Cardboard triangle balancing on a fingertip. ORGANIC (pictorial) -- traced
% against scan02 p12 (book p.16). The medians are exact; the hand is stylised.
\newcommand{\FIGITWENTYTHREE}{%
\begin{bkfigure}[0.50\linewidth]
\begin{tikzpicture}[scale=1.0]
  % triangle seen almost edge-on, balanced at its centroid
  \coordinate (A) at (-1.55,0.18);
  \coordinate (B) at (1.62,0.30);
  \coordinate (C) at (-0.28,0.92);
  \coordinate (Ma) at ($(B)!0.5!(C)$);
  \coordinate (Mb) at ($(C)!0.5!(A)$);
  \coordinate (Mc) at ($(A)!0.5!(B)$);
  \coordinate (O)  at ($(A)!0.6666666!(Ma)$);
  \draw[given] (A) -- (B) -- (C) -- cycle;
  \draw[aux] (A) -- (Ma);
  \draw[aux] (B) -- (Mb);
  \draw[aux] (C) -- (Mc);
  % stylised finger and hand: the rounded fingertip touches the centroid and
  % the finger runs down to meet the top of the fist (it must not float free)
  \draw[fig] ($(O)+(-0.075,-1.14)$) -- ($(O)+(-0.075,-0.02)$)
             arc (180:0:0.075) -- ($(O)+(0.075,-1.14)$);
  \draw[fig] (-0.30,-0.62) .. controls (-0.34,-0.98) and (-0.30,-1.24) .. (-0.24,-1.42);
  \draw[fig] ( 0.16,-0.62) .. controls (0.22,-0.98) and (0.24,-1.24) .. (0.20,-1.42);
  \draw[fig] (-0.30,-0.62) .. controls (-0.20,-0.70) and (0.06,-0.70) .. (0.16,-0.62);
  \foreach \k in {0,1,2} {
    \draw[fig] (-0.30,-0.78-0.17*\k)
       .. controls (-0.40,-0.84-0.17*\k) and (-0.42,-0.92-0.17*\k) ..
          (-0.32,-0.96-0.17*\k);}
  \draw[fig] (-0.30,-1.42) rectangle (0.22,-1.60);
\end{tikzpicture}
\figcap{1--23}
\end{bkfigure}}

%--------------------------------------------------------------- Fig 1-24
% Altitudes. (a) acute, foot inside; (b) obtuse, foot on the extension of the
% base; (c) the three altitude lines concurrent at O. All feet are exact
% projections ($(X)!(P)!(Y)$).
\newcommand{\FIGITWENTYFOUR}{%
\begin{bkfigure}[\linewidth]
\begin{tikzpicture}[scale=1.0]
  % Proportions measured off the book: all three panels are about the same width
  % (100 : 106 : 110) and grow in height (64 : 92 : 113).
  % ---- (a) apex left of centre, foot inside, right-angle box right of the altitude
  \begin{scope}
    \coordinate (A) at (0,0);
    \coordinate (B) at (2.15,0);
    \coordinate (C) at (0.7256,1.376);
    \coordinate (F) at ($(A)!(C)!(B)$);
    \draw[given] (A) -- (B) -- (C) -- cycle;
    \draw[aux] (C) -- (F);
    \draw[fig] ($(F)+(0.16,0)$) -- ($(F)+(0.16,0.16)$) -- ($(F)+(0,0.16)$);
    \node[slab, below=6pt] at (1.075,0) {(a)};
  \end{scope}
  % ---- (b) obtuse at B: the foot falls beyond B
  \begin{scope}[xshift=3.05cm]
    \coordinate (A) at (0,0);
    \coordinate (B) at (1.3975,0);
    \coordinate (C) at (2.2844,1.9834);
    \coordinate (F) at ($(A)!(C)!(B)$);
    \draw[given] (A) -- (B) -- (C) -- cycle;
    \draw[given] (B) -- (F);
    \draw[aux] (C) -- (F);
    \draw[fig] ($(F)+(-0.16,0)$) -- ($(F)+(-0.16,0.16)$) -- ($(F)+(0,0.16)$);
    \node[slab, below=6pt] at (1.142,0) {(b)};
  \end{scope}
  % ---- (c) three altitudes concurrent; near-isosceles, apex at 0.584 of the base
  \begin{scope}[xshift=6.10cm]
    \coordinate (A) at (0,0);
    \coordinate (B) at (2.365,0);
    \coordinate (C) at (1.3812,2.4295);
    \coordinate (Fc) at ($(A)!(C)!(B)$);
    \coordinate (Fa) at ($(B)!(A)!(C)$);
    \coordinate (Fb) at ($(A)!(B)!(C)$);
    \coordinate (H)  at (intersection of C--Fc and A--Fa);
    \draw[given] (A) -- (B) -- (C) -- cycle;
    \draw[aux] (C) -- (Fc);
    \draw[aux] (A) -- (Fa);
    \draw[aux] (B) -- (Fb);
    \draw[fig] ($(H)+(-0.08,-0.08)$) -- ($(H)+(0.08,0.08)$);
    \draw[fig] ($(H)+(-0.08,0.08)$) -- ($(H)+(0.08,-0.08)$);
    % 176.21 deg bisects the gap between the altitude rays at 150.38 and 202.05
    \node[vlab] at ($(H)+(176.21:0.65)$) {$O$};
    \node[slab, below=6pt] at (1.1825,0) {(c)};
  \end{scope}
\end{tikzpicture}
\figcap{1--24}
\end{bkfigure}}

%--------------------------------------------------------------- Fig 1-25
% The three angle bisectors meet at the incentre. Each bisector is drawn from a
% vertex through the exact incentre to the opposite side.
\newcommand{\FIGITWENTYFIVE}{%
\begin{bkfigure}[0.52\linewidth]
\begin{tikzpicture}[scale=1.0]
  % proportions from the book: apex at 0.775 of the base, height 0.667 of the base
  \coordinate (A) at (0,0);
  \coordinate (B) at (2.95,0);
  \coordinate (C) at (2.28625,1.96765);
  % incentre = (a A + b B + c C)/(a+b+c), a=|BC|, b=|CA|, c=|AB|
  % |BC| = 2.077003, |CA| = 3.016421, |AB| = 2.95  ->  I = (1.944900, 0.721694)
  \coordinate (I) at (1.944900,0.721694);
  \draw[given] (A) -- (B) -- (C) -- cycle;
  \draw[aux] (A) -- (intersection of A--I and B--C);
  \draw[aux] (B) -- (intersection of B--I and C--A);
  \draw[aux] (C) -- (intersection of C--I and A--B);
  \draw[fig] ($(I)+(-0.08,-0.08)$) -- ($(I)+(0.08,0.08)$);
  \draw[fig] ($(I)+(-0.08,0.08)$) -- ($(I)+(0.08,-0.08)$);
  % 227.52 deg bisects the gap between the bisector rays at 200.36 and 254.68
  \node[vlab] at ($(I)+(227.52:0.58)$) {$O$};
\end{tikzpicture}
\figcap{1--25}
\end{bkfigure}}

%--------------------------------------------------------------- Figs 1-26, 1-27
% 1-26: l parallel to m (same direction vector), transversal n; angles 1 and 2
%       are corresponding angles, hence equal.
% 1-27: copying angle 2 at P produces a line through P parallel to l.
\newcommand{\FIGITWENTYSIXTWENTYSEVEN}{%
\begin{bkfigure}[\linewidth]
\begin{minipage}[t]{0.48\linewidth}\centering
\begin{tikzpicture}[scale=1.0]
  % l and m share the direction (3.2, 0.9); n is the transversal
  \coordinate (l1) at (0.10,0.95);  \coordinate (l2) at (3.30,1.85);
  \coordinate (m1) at (0.10,0.30);  \coordinate (m2) at (3.30,1.20);
  \coordinate (n1) at (0.55,2.05);  \coordinate (n2) at (2.05,-0.10);
  \draw[given] (l1) -- (l2);
  \draw[given] (m1) -- (m2);
  \draw[given] (n1) -- (n2);
  \coordinate (P1) at (intersection of l1--l2 and n1--n2);
  \coordinate (P2) at (intersection of m1--m2 and n1--n2);
  % the book marks the two ALTERNATE INTERIOR angles: 1 below l and right of n,
  % 2 above m and left of n. Each label sits on the bisector of its own sector
  % (l runs at 15.71 deg, n descends at 304.90 deg).
  \node[slab] at ($(P1)+(340.30:0.42)$) {1};
  \node[slab] at ($(P2)+(160.30:0.42)$) {2};
  \node[vlab, right] at (l2) {$l$};
  \node[vlab, right] at (m2) {$m$};
  \node[vlab, below right, xshift=-1pt, yshift=-0pt] at (n2) {$n$};
\end{tikzpicture}
\figcap{1--26}
\end{minipage}\hfill
\begin{minipage}[t]{0.48\linewidth}\centering
\begin{tikzpicture}[scale=1.0]
  \coordinate (l1) at (0,1.35);  \coordinate (l2) at (3.10,1.35);
  \coordinate (n1) at (0.72,0);  \coordinate (n2) at (2.02,2.10);
  \draw[given] (l1) -- (l2);
  \draw[given] (n1) -- (n2);
  \coordinate (Q) at (intersection of l1--l2 and n1--n2);
  \coordinate (P) at ($(n1)!0.30!(n2)$);
  \draw[aux] ($(P)+(-0.66,0)$) -- ($(P)+(1.22,0)$);   % the copied parallel
  \fill (P) circle (1.2pt);
  % as in the book: 2 is the lower-left angle at Q, 1 the upper-right angle at P
  % (alternate interior angles). n rises at 58.24 deg.
  \node[slab] at ($(Q)+(209.12:0.55)$) {2};
  \node[slab] at ($(P)+(29.12:0.55)$) {1};
  \node[vlab] at ($(P)+(299.12:0.34)$) {$P$};
  \node[vlab, left] at (l1) {$l$};
  \node[vlab, above] at (n2) {$n$};
\end{tikzpicture}
\figcap{1--27}
\end{minipage}
\end{bkfigure}}

%--------------------------------------------------------------- Fig 1-28
% Right triangle: right angle at C, hypotenuse c, legs a and b.
\newcommand{\FIGITWENTYEIGHT}{%
\begin{bkfigure}[0.52\linewidth]
\begin{tikzpicture}[scale=1.0]
  \coordinate (A) at (0,0);
  \coordinate (C) at (2.45,0);
  \coordinate (B) at (2.45,1.70);
  \draw[given] (A) -- (C) -- (B) -- cycle;
  \draw[fig] (2.29,0) -- (2.29,0.16) -- (2.45,0.16);
  \node[vlab, left]  at (A) {$A$};
  \node[vlab, below right, xshift=-1pt, yshift=-0pt] at (C) {$C$};
  \node[vlab, above] at (B) {$B$};
  \node[vlab, below] at (1.22,0) {$b$};
  \node[vlab, right] at (2.45,0.85) {$a$};
  \node[vlab, above left, xshift=-0pt, yshift=-1pt] at (1.22,0.85) {$c$};
\end{tikzpicture}
\figcap{1--28}
\end{bkfigure}}
